\section*{第 \ref{ch:g2:caring-for-nature} 章——关爱自然}

\begin{solution}{exo:g2:caring-for-nature:1}
\omterm{def:g1:living-or-not:living}{生物}赖以生存的外界。例如：我们呼吸的空气、
我们喝的水、我们吃的食物——都来自
\omterm{def:g2:caring-for-nature:environment}{自然}。
\end{solution}

\begin{solution}{exo:g2:caring-for-nature:2}
例如：树皮下的甲虫、
腐木下的鼠妇、表面的苔藓，以及在树洞里过冬的刺猬。
\end{solution}

\begin{solution}{exo:g2:caring-for-nature:3}
例如：把所有垃圾带回家；只看不采，
不要把一丛\omterm{def:g1:plants-around-us:plant}{植物}全部摘光；翻开的石头或
木头要放回原处；在易受损的地方要沿小路行走；
远远地安静观察\omterm{def:g1:animals-around-us:animal}{动物}，不要打扰它们的幼体。
\end{solution}

\begin{solution}{exo:g2:caring-for-nature:4}
因为大\omterm{def:g2:caring-for-nature:environment}{自然}没有垃圾车：我们留下的东西会一直留在那里——
一个塑料盖子可能要好多好多年，足以困住或毒害
\omterm{def:g1:animals-around-us:animal}{动物}，等今天的孩子长大成人时，
它还躺在那里。
\end{solution}

\begin{solution}{exo:g2:caring-for-nature:5}
这一丛若被采光，就无法结出\omterm{def:g2:from-seed-to-plant:seed}{种子}——
明年这里便不会再有雏菊，蜜蜂也失去了取食的花。
\end{solution}

\begin{solution}{exo:g2:caring-for-nature:6}
它们会慢慢腐烂，变成深色、松软的土壤。
蚯蚓、鼠妇和微小的\omterm{def:g1:living-or-not:living}{生物}都在其中出力；最终我们得到
堆肥——一种能为来年\omterm{def:g1:plants-around-us:plant}{植物}提供养分的好土壤。
\end{solution}

\begin{solution}{exo:g2:caring-for-nature:7}
例如：热浪来袭时放一碟水，
留一角不修剪的野草，天冷时为鸟类准备\omterm{def:g2:from-seed-to-plant:seed}{种子}，
挖一个小池塘。
\end{solution}

\begin{solution}{exo:g2:caring-for-nature:8}
答案不唯一。塑料包装袋和瓶子通常最多。
每一种垃圾都应该跟随主人带回家，
再放进正确的分类垃圾桶。
\end{solution}

\begin{solution}{exo:g2:caring-for-nature:9}
让\omterm{ex:g2:animals-grow-up:frog}{蝌蚪}留在池塘里：池塘里有它们需要的一切——
池水、食物、长出腿的空间——而广口瓶里没有。如果
艾玛想帮忙，她可以在春季里在那里观察它们，
保持岸边不被踩踏，或许还能帮忙建一个像学校那样的新池塘\omterm{def:g2:where-animals-live:habitat}{栖息地}。
保持岸边不被踩踏，或许还能帮忙建一个像学校那样的新池塘\omterm{def:g2:where-animals-live:habitat}{栖息地}。
\end{solution}

\begin{solution}{exo:g2:caring-for-nature:10}
向坏的方向，一百万个“无所谓”就会让一片草地被瓶盖埋没——
小伤害会成倍增加。向好的方向，一百万份
微小的用心——把垃圾带回家、随手关水龙头、把废弃物分类——
同样会积少成多，
大\omterm{def:g2:caring-for-nature:environment}{自然}大多是靠这样细小的选择来守护的。
\end{solution}
